\subsection{Agency statistics and the one-year default rate}
\label{sec:rw-agency}

The published statistics on which a validator relies come from the agencies themselves and, for the European entities, from ESMA's CEREP repository, whose help file describes a stock model: intra-period rating activity is derived from a comparison of the ratings at the beginning and at the end of a period \citep{ESMA_CEREP}. The supervisory definition of the one-year default rate is that of the guidelines of the European Banking Authority: the denominator holds the non-defaulted obligors observed at the beginning of the year, the numerator those among them with at least one default event during the year, an obligor whose obligation was closed stays in both, and each defaulted obligor is counted once \citep{EBA_2017}. Schuermann and Hanson describe the opposite convention as the one that has emerged as an industry standard, in which withdrawn ratings are gradually eliminated from the sample and the probability of withdrawal is distributed among all states in proportion to their values \citep{Schuermann_2004}. Cantor and Hamilton address the adjustment of corporate default rates for rating withdrawals \citep{Cantor_2007}, and Carey and Hrycay study how credit risk models are parameterised with rating data \citep{Carey_2000}. The Basel Committee's studies of the validation of internal rating systems frame the comparison of such estimates with observed default rates as a task of validation \citep{BCBS_2005}. None of these works measures, on data whose truth is known, how far each convention moves the PD of a grade when withdrawal is informative, and none translates the difference between conventions into risk weight.

\subsection{Estimators of migration and default probabilities}
\label{sec:rw-estimators}

The cohort estimator divides the number of transitions from one grade to another by the number of obligors in the grade at the start of the year; Anderson and Goodman derived its maximum-likelihood properties for a Markov chain \citep{Anderson_1957}. The duration estimator of Lando and Skodeberg uses the time spent in each state and the dates of the changes, and so assigns a positive PD to a grade in which the cohort estimate is zero \citep{Lando_2002,Schuermann_2004}; the idea of weighting each observation by the time it was at risk is that of the product-limit estimator of Kaplan and Meier for a single event \citep{Kaplan_1958}. Jafry and Schuermann compare migration matrices and summarise their mobility \citep{Jafry_2004}. When only the annual counts are published, as in CEREP, a generator of the continuous-time chain can be recovered from the matrix by the adjustments of Israel, Rosenthal and Wei \citep{Israel_2001}, by the approximation of Jarrow, Lando and Turnbull \citep{Jarrow_1997}, or by the expectation-maximisation algorithm for a discretely observed jump process of Bladt and S\o rensen \citep{Bladt_2005} as presented for credit migrations by dos Reis and Smith \citep{dosReis_2018}. The Markov and time-homogeneity assumptions have been relaxed by mixtures of Markov chains \citep{Frydman_2008}, by dependence on the business cycle \citep{Bangia_2002} and by a self-exciting downgrade intensity that models rating momentum \citep{dosReis_2020}. All of these estimators consume the withdrawn ratings through a convention chosen before estimation; none of them changes it.

\subsection{Informative withdrawal as informative censoring and missing data}
\label{sec:rw-censoring}

A withdrawn rating is a right-censored observation of the time to default, and the censoring is informative when its hazard depends on the proximity of default. The product-limit estimator assumes independent censoring \citep{Kaplan_1958}. Tsiatis showed that the joint distribution of latent failure times in a competing-risks model is not identifiable from the observed data \citep{Tsiatis_1975}, so that no analysis of the withdrawals alone can separate informative from uninformative withdrawal. Cause-specific hazards \citep{Prentice_1978} and the subdistribution hazard of a competing risk \citep{Fine_1999} are the standard ways to estimate the incidence of an event in the presence of another that prevents it, and each answers a stated question. Rubin's classification of missing-data mechanisms separates the cases in which the mechanism can be ignored from those in which it cannot \citep{Rubin_1976,Little_2019}. For the outcomes that are missing not at random, Manski's worst-case bounds give the identified set without any assumption on the mechanism \citep{Manski_1989}, and the tipping-point analysis of clinical trials varies a sensitivity parameter until a conclusion changes \citep{Cro_2020}. The protocol uses the missing-data view to define what is measured, the PD of a grade over a year for the ratings that would have stayed outstanding, and uses the generator of paths to produce the truth that an observer of real data never has.

\subsection{Intervals for the PD of a thin grade}
\label{sec:rw-intervals}

The Wald interval around an observed default rate is the one that most reports use, and Brown, Cai and DasGupta show that its coverage of a binomial proportion is erratic, with a systematic shortfall for small proportions, and recommend the Wilson, Jeffreys and Agresti-Coull intervals \citep{Brown_2001}. The Wilson interval inverts the score test \citep{Wilson_1927}, the Agresti-Coull interval adds pseudo-observations to the Wald form \citep{Agresti_1998}, the Clopper-Pearson interval inverts the binomial test exactly \citep{Clopper_1934}, and Newcombe compares seven methods on coverage and length \citep{Newcombe_1998}. For a count of zero, Hanley and Lippman-Hand derive the upper bound that applies when nothing goes wrong \citep{Hanley_1983}. Hanson and Schuermann compare analytical intervals, among them the Wald interval, with bootstrapped intervals for cohort and duration estimates of the PDs of rating grades; they find the bootstrapped intervals of the duration estimates tighter than the Wald interval, and find that even these tight intervals cannot tell notch-level PDs of the investment grades apart, while those of the speculative grades separate cleanly \citep{Hanson_2006,Schuermann_2004}. The calibration of PD curves, which carries a PD observed in one period to another, takes the observed default rate by grade as its input and so inherits its definition \citep{Tasche_2013}. The gap the protocol fills is the exact coverage by enumeration at the true PDs of an agency-like generator, grade by grade, with the closed form of the shortfall.

\subsection{Evaluation on known truth and the translation to capital}
\label{sec:rw-evaluation}

Morris, White and Crowther recommend that a simulation study state its aims, its data-generating mechanism, the estimands, the methods compared and the performance measures, and report the Monte Carlo standard error of each \citep{Morris_2019}; the protocol follows that structure. The capital implied by a PD follows from the one-factor model of Gordy, which underlies the IRB risk-weight function \citep{Gordy_2003,BCBS_CRE31}. Engelmann shows that projections of a default rate from a migration matrix drift towards a portfolio of the matrix's own, its through-the-cycle portfolio, even when no scenario is applied, an effect of the same family as those measured here, in which an estimator's convention changes the PD with no change in the obligors \citep{Engelmann_2024}. The bootstrap over cohort years, which provides the interval of the definition gap, resamples units of an annual series \citep{Efron_1979}; a block scheme would be needed to respect a dependence between years \citep{Kunsch_1989}, and Section~\ref{sec:limitations} states the consequence for the interval of the definition gap.

\subsection{Where the proposed protocol stands}
\label{sec:rw-position}

The works above fall into four groups: the conventions for counting defaults and withdrawals \citep{EBA_2017,Schuermann_2004,Cantor_2007,ESMA_CEREP}, the estimators that consume them \citep{Anderson_1957,Lando_2002,Israel_2001,Bladt_2005}, the statistical theory of informative censoring and missing data \citep{Tsiatis_1975,Rubin_1976,Manski_1989}, and the intervals for a binomial proportion \citep{Brown_2001,Wilson_1927,Agresti_1998}. The proposed protocol stands between them. It takes the conventions of the first group as parameters, runs the estimators of the second on data generated by a chain whose truth is known, uses the missing-data view of the third to define the bias it measures, and computes the coverage of the fourth exactly at the PDs of rating grades. Its novelty is the combination and the measurement, not a new estimator: it quantifies, for the conventions the supervisory text and the agencies' pages offer, the size of the PD error that a validator accepts by choosing one, and the size of the capital error that follows from it. The measurement is reproducible with the published engine and the agencies' public statistics, which the works above do not provide together.
